% =========================================================
% Why Induction Works
% =========================================================
\paragraph{Why Induction Works.}

Mathematical induction is often presented as a technique: base case,
inductive step, done. That presentation is useful but incomplete.
Induction \emph{feels} like a procedure, but it is actually a theorem,
and it is worth understanding why it is true before using it.

\begin{remark}[The Peano axiom behind induction]
The natural numbers are characterized by the Peano axioms.
The critical one for induction is Axiom~P5:

\medskip
\noindent\textit{If $A \subseteq \mathbb{N}$ satisfies (i) $0 \in A$ and
(ii) $n \in A \Rightarrow S(n) \in A$, then $A = \mathbb{N}$.}

\medskip
\noindent This axiom says that $\mathbb{N}$ is the \emph{smallest}
set containing $0$ and closed under the successor function $S$.
There is no proper subset of $\mathbb{N}$ with these two properties.

The induction principle is a direct translation. Let $P(n)$ be any
statement. Define $A = \{n \in \mathbb{N} : P(n)\}$. If $P(0)$ holds,
then $0 \in A$. If $P(n) \Rightarrow P(S(n))$ for every $n$, then
$A$ is closed under $S$. By P5, $A = \mathbb{N}$.
Therefore $P(n)$ holds for every $n \in \mathbb{N}$.

The full development of this from the Peano axioms is in the
Axiom Systems chapter. Here, the takeaway is: induction is not
a trick. It is a theorem derivable from the definition of $\mathbb{N}$.
\end{remark}

\begin{remark}[Why the base case is not optional]
The inductive step alone says: \emph{if} $P(n)$ holds for some $n$,
then $P(n+1)$ holds. This is a conditional, not an assertion.
Without a base case, there is no $n$ for which $P(n)$ is known to hold.
The chain of conditionals has no starting point, and the proof produces
nothing.

The most famous illustration is the false theorem ``all horses are the
same colour.'' A flawed argument runs: base case trivial ($n=1$),
inductive step [fatally flawed reasoning]. The lesson is not just that
the step is wrong --- it is that even a correct base case does not rescue
a bad step, and even a correct step does not rescue a missing base case.
Each part is independent; both are necessary.

In terms of P5: the base case establishes $0 \in A$. The inductive step
establishes closure under $S$. Both conditions of P5 are required for
the conclusion $A = \mathbb{N}$.
\end{remark}

\begin{remark}[The domino picture and its limitations]
The common picture of induction as an infinite row of falling dominoes
is a useful mnemonic but a misleading explanation. It suggests that
induction requires each domino to knock over the next, which captures
the inductive step. But it does not explain why the argument is valid ---
the validity comes from P5, not from a physical metaphor.

More importantly, the domino picture suggests that the steps happen
\emph{in time}: first domino~0 falls, then~1, then~2, and so on.
But induction proves a statement about \emph{all} natural numbers
simultaneously, not by a sequential process. The proof is instantaneous:
once you have established the base case and the inductive step, the
conclusion $P(n)$ for all $n$ follows at once from P5.

Use the domino picture as memory aid. For understanding, use P5.
\end{remark}

\begin{remark}[Induction and well-ordering are equivalent]
The induction principle and the well-ordering principle for $\mathbb{N}$
(every nonempty subset has a least element) are logically equivalent.
This means every induction argument can be rephrased as a
minimal-counterexample argument, and vice versa. Both are valid proof
tools. The well-ordering approach is sometimes more natural when you
want to reason about the \emph{first} failure rather than climbing
from a base case. The equivalence is proved in
Section~\ref{sec:wellordering}.
\end{remark}
